\title{Availability, reproducibility and reusability evaluation levels: guideline for authors}

\section{Data management}


\section{Code management}
We distinguished 3 different independent categories (some badges may be obtained independently):

\begin{itemize}
\item Availability: the software artefacts used in the publication adhere to the standard Open Science standard. Reader can find the code in an open repository
\item Reproducibility: readers can reproduce the main results of the publication. The authors have provided all the necessary elements to enable this.
  
\item Reusability: The software artefacts can be reused for goals or experiments other than those described in the article. The authors have provided readers with all the elements needed to modify the code, add new functionalities, and ensure cumulative science. 
\end{itemize}

For each of them, a finer granularity is defined to indicate what can be done to improve their work.


\subsection{Availability of the code}

The code used in the publication has to be available for reviewers and readers. There are various level of availability of the code detailed in the following table:\\


\begin{tabular}{|p{10cm}|p{1.5cm}|}
%\begin{tabular}{|c|c|}
  \hline
  \cellcolor{bl1}  \cellcolor{bl1} Description & \cellcolor{bl1} Stamp \\
\hline
code available on a public forge + license, authors, contributors files, commit number and branche corresponding to the publication & \textcolor{green1}{\ding{78}} \\
\hline
+ code available on Zenodo & \textcolor{green1}{\ding{78}\ding{78}} \\
\hline
+ Software Heritage Identifier (SWHID) &  \textcolor{green1}{\ding{78}\ding{78}\ding{78}}\\

\hline
+ HAL deposit & \textcolor{green1}{\ding{78}\ding{78}\ding{78}\ding{78}} \\ 
\hline
\end{tabular}
~\\

More details :
\begin{itemize}
\item Level 1: \\
without a licence file, the code cannot be used by other people. The files of the authors and contributors are needed for patrimonial and author's rights. Finally, providing the commit and the exact branch enables further modifications to be made to the repository without losing the state corresponding to the article.   

\item Level 2: \\
Zenodo provides a DOI and enables to fix a precise state of the software
  
\item Level 3: \\
Hosting your source code on the Software Heritage platform enables a perennial archive and a very precise commit identification (SWHID) 
\item Level 4: \\
  HAL deposit enables to verify that all the needed files are present and in the right format (curation)\\
\end{itemize}

%%%%%%%%%%%%%%%%%%%%%%%%%%%%%%%%%%%%%%%%%%

\subsection{Reproducibility}

\begin{tabular}{|p{10cm}|p{1cm}|p{1.7cm}|}
\hline
 \cellcolor{bl1} Description & \cellcolor{bl1} Stamp & \cellcolor{bl1} Requirement \\
\hline
re-runable code, description of the workflow and test case(s) &  \textcolor{yell1}{\ding{115}} & \textcolor{green1}{\ding{78}}\\ 
\hline
+ workflow manager for automatic execution  &  \textcolor{yell1}{\ding{115}\ding{115}} & \textcolor{green1}{\ding{78}}\\ 
\hline
+ reproducible software environment & \textcolor{yell1}{\ding{115}\ding{115}\ding{115}} & \textcolor{green1}{\ding{78}} \\
\hline
+ regression, test cases and documentation in CI & \textcolor{yell1}{\ding{115}\ding{115}\ding{115}\ding{115}} &  \textcolor{green1}{\ding{78}} \\
 \hline
\end{tabular}
~\\

More details :
\begin{itemize}
\item Level 1: \\
Supposing the code is available and usable, the reader can re-run understanding with component has to be used and in which order. He can also verify througgh a test that he obtains correct results.    

\item Level 2: \\
A workflow manager solution is provided: the reader has just to launch the procedure in order to run a test case.
  
\item Level 3: \\
The software environement is described: the reader can deploy it in his own computer autoatically or at leat manually with the correct version number of the dependances

\item Level 4: \\
The repository contains regression tests, test cases and automatic documentation in a CI: any modification is tested and does not modify the behaviour of the code. Moreover, new documentation is automatically generated.   
\end{itemize}



%%%%%%%%%%%%%%%%%%%%%%%%%%%%%%%%%%%%%%%%%%

\subsection{Reusability for further projects}


\begin{tabular}{|p{10cm}|p{1cm}|p{1.7cm}|}
\hline
\cellcolor{bl1} Description & \cellcolor{bl1} Stamp & \cellcolor{bl1} Requirement \\
\hline
All functions documentation inside the code & \textcolor{purple1}{\ding{112}} &  \textcolor{yell1}{\ding{115}} \\
+ Description of each element in a Readme + workflow description for execution &  \textcolor{purple1}{\ding{112}\ding{112}} &\\
+ User documentation with examples &  \textcolor{purple1}{\ding{112}\ding{112}\ding{112}}&\\
+ Developper documentation (``How to contribute ?'' guide) &  \textcolor{purple1}{\ding{112}\ding{112}\ding{112}\ding{112}}&\\
\hline
\end{tabular}
~\\

More details :
\begin{itemize}
\item Level 1: \\
Each module of the code is highly documented with useful tools as sphynx or doxygen for example  

\item Level 2: \\
The role of each piece of code or script is documented in a Readme file. Each step of execution is described in order to facilitate the reusability.   
\item Level 3: \\
The ``user'' documentation is complete with simple test cases, examples, links with bibliography and theoretical elements. Numerical methods are described, as is the rationale behind their selection.   
\item Level 4: \\
The repository contains a guide for future developper, explaining the choices made during the software development and the way to modify it. 
\end{itemize}





%%% Local Variables:
%%% mode: latex
%%% TeX-master: "main"
%%% End:
